\begin{examplebox}
\textbf{例 1：由方程求焦点、准线。}
$x^2/25+y^2/9=1$ 中 $a=5,c=4$，故右焦点为 $\boxed{(4,0)}$，右准线为 $\boxed{x=25/4}$。
回验：$P=(0,3)$ 时 $PF_R=5$，准距为 $25/4$，比值为 $4/5$。

\medskip
\textbf{例 2：求准距应除以 $e$。}
$x^2/16+y^2/12=1$ 上，若 $PF_R=4$，则 $a=4,c=2,e=1/2$，
故 $d(P,\ell_R)=PF_R/e=\boxed{8}$。
回验：由 $4=4-x/2$ 得 $x=0$，$P=(0,\pm2\sqrt3)$，距离到 $(2,0)$ 确为 $4$，到准线 $x=8$ 为 $8$。
\end{examplebox}
\clearpage
\begin{examplebox}
\textbf{例 3：系数恰好为 $e$，先统一距离。}
在 $x^2/25+y^2/9=1$ 上求 $PF_R+\tfrac45d(P,\ell_R)$ 的最小值及取等点。
\begin{center}\begin{tikzpicture}[font=\small,>=stealth]
\draw[black!30,rounded corners=2pt] (0,0) rectangle (14.4,8.2);
\node[anchor=west,font=\small\bfseries] at (.5,7.72) {例 3：右顶点最靠近右准线，焦半径也最短};
\begin{scope}[shift={(6.3,3.8)},scale=.58]
\draw[blue,very thick] (0,0) ellipse (5 and 3);
\draw[violet,thick] (6.25,-3.5)--(6.25,4.3) node[above] {$\ell_R$};
\draw[black!40] (-5.3,0)--(6.6,0);
\draw[orange!85!black,very thick] (5,0)--(6.25,0);
\draw[violet,dashed,thick] (0,3)--(6.25,3);
\fill (5,0) circle (2pt);
\node[anchor=west] at (7.8,-1.2) {$A(5,0)$};
\draw[black!45,->] (7.6,-.9)--(5,0);
\fill (0,3) circle (2pt) node[above left] {$B(0,3)$};
\fill (4,0) circle (2pt) node[below left=7pt] {$F_R$};
\node[above] at (3.0,3) {$d(B,\ell_R)=25/4$};
\end{scope}
\node at (7.2,1.2) {$d(A,\ell_R)=5/4,\quad AF_R=1$};

\node[anchor=west] at (.7,.55) {$PF_R+\frac45d(P,\ell_R)=2PF_R=10-\frac85x$；$x=5$ 时最小值为 $2$};
\end{tikzpicture}\end{center}

由定义，原式 $=2PF_R=\tfrac85d(P,\ell_R)$。
又 $x\in[-5,5]$，$d(P,\ell_R)=25/4-x\ge5/4$，
所以原式 $\ge\tfrac85\cdot\tfrac54=\boxed{2}$。
仅当 $x=5$ 时取等，椭圆方程给出 $y=0$，故唯一取等点为 $\boxed{(5,0)}$。
回验：此点焦半径为 $1$、准距为 $5/4$，原式确为 $1+1=2$。
\end{examplebox}
\clearpage

\begin{examplebox}
\textbf{例 4：反向应用——给焦点、准线、比例，写轨迹。}
固定焦点 $F=(2,0)$、准线 $\ell:x=8$，求满足
$PF=\tfrac12 d(P,\ell)$ 的所有点的轨迹。
\textbf{不必先设椭圆半轴，直接用距离公式：}
\[
\begin{aligned}
\sqrt{(x-2)^2+y^2}&=\frac12|x-8|\\
\Longleftrightarrow\quad (x-2)^2+y^2&=\frac14(x-8)^2\\
\Longleftrightarrow\quad \frac34 x^2+y^2&=12\\
\Longleftrightarrow\quad \boxed{\frac{x^2}{16}+\frac{y^2}{12}=1}.
\end{aligned}
\]
\textbf{核验：} 新椭圆 $a=4,c=2,e=1/2$，右准线 $x=a^2/c=8$，与题设完全一致。
由于平方等价，两边都是非负长度，所以没有引入增根。
\end{examplebox}
\clearpage
\begin{examplebox}
\textbf{例 5：换成 $y$ 轴——不是只有竖直椭圆图形，准线也要转向。}
椭圆 $\frac{y^2}{25}+\frac{x^2}{16}=1$ 的上焦点为
$F_U=(0,3)$，上准线为 $\ell_U:y=25/3$，离心率 $e=3/5$。
取椭圆上的点 $P=(4,0)$，则
\[
PF_U=\sqrt{4^2+(-3)^2}=5,\qquad
 d(P,\ell_U)=\frac{25}{3},\qquad
 \frac{PF_U}{d(P,\ell_U)}=\frac35=e.
\]
\textbf{口头规则：} 长轴沿 $x$，准线写 $x=\pm a/e$；长轴沿 $y$，准线写 $y=\pm a/e$。
\end{examplebox}
\clearpage
\begin{examplebox}
\textbf{例 6：只记住“焦点坐标 + 离心率”，快速找准线。}
某椭圆中心在原点，长轴在 $x$ 轴上，右焦点为 $(3,0)$，离心率为 $e=3/5$。
求右准线和标准方程。
\[
a=\frac ce=\frac{3}{3/5}=5,\qquad
\boxed{\ell_R:x=\frac ae=\frac{25}{3}},\qquad
b^2=a^2-c^2=25-9=16.
\]
所以椭圆为 $\boxed{x^2/25+y^2/16=1}$。
\textbf{快捷记法：} 焦点 $c$ 已知，先“$c\div e$ 得 $a$”，再“$a\div e$ 得准线位置”。
\end{examplebox}
